\documentclass[a4paper, 11pt]{article}

%% ------------------------------------------------------------
%%  PAQUETES Y CONFIGURACIÓN BÁSICA
%% ------------------------------------------------------------
\usepackage[utf8]{inputenc}
\usepackage[T1]{fontenc}
\usepackage[spanish, es-tabla]{babel}
\usepackage{lmodern}
\usepackage[top=2.0cm, bottom=2.0cm, left=2.3cm, right=2.3cm]{geometry}
\usepackage{graphicx}
\graphicspath{{Figures/}}
\usepackage{setspace}
\setstretch{1.08}
\setlength{\parskip}{0.28em}
\setlength{\parindent}{0pt}

\usepackage{enumitem}
\setlist{noitemsep, topsep=1.5pt, parsep=0pt, partopsep=0pt}

\usepackage[dvipsnames]{xcolor}
\definecolor{uni-blue}{RGB}{0, 56, 101}

\usepackage[
    colorlinks = true,
    linkcolor  = uni-blue,
    citecolor  = uni-blue,
    urlcolor   = uni-blue
]{hyperref}

\usepackage{fancyhdr}
\usepackage{titlesec}

%% Estilo de títulos compacto y elegante
\titleformat{\section}
    {\normalfont\large\bfseries\color{uni-blue}}{\thesection.}{0.5em}{}
    [\vspace{1pt}\titlerule\vspace{2pt}]

\titlespacing*{\section}{0pt}{0.65em}{0.2em}

%% Encabezados y pie de página
\pagestyle{fancy}
\fancyhf{}
\setlength{\headheight}{14.5pt}
\fancyhead[L]{\footnotesize\textit{Ética y Filosofía Política --- Sección B}}
\fancyhead[R]{\footnotesize\textit{I. Kant: ¿Por qué la metafísica no es ciencia?}}
\fancyfoot[C]{\small\thepage}
\renewcommand{\headrulewidth}{0.4pt}

%% ------------------------------------------------------------
%%  DATOS DEL TRABAJO
%% ------------------------------------------------------------
\newcommand{\univname}{UNIVERSIDAD NACIONAL DE INGENIERÍA}
\newcommand{\facultad}{FACULTAD DE CIENCIAS}
\newcommand{\escuela}{Escuela Profesional de Física}
\newcommand{\curso}{Ética y Filosofía Política --- Sección B}
\newcommand{\docente}{Prof. José Romero Cieza}
\newcommand{\alumno}{Vera Vivanco, Jesús Rodolfo}
\newcommand{\codigo}{20252703K}
\newcommand{\fecha}{Lima, 30 de octubre de 2026}
\newcommand{\tema}{¿Por qué la metafísica no puede llegar a ser una ciencia?}
\newcommand{\subtema}{Análisis a partir del Prólogo a la segunda edición (1787) de la \textit{Crítica de la razón pura} de Immanuel Kant}

\begin{document}

%% ------------------------------------------------------------
%%  PORTADA OFICIAL (PÁGINA 1)
%% ------------------------------------------------------------
\begin{titlepage}
    \centering
    \vspace*{0.6cm}
    {\Large\bfseries \univname}\\[0.25cm]
    {\large\bfseries \facultad}\\[0.15cm]
    {\normalsize \escuela}\\[0.9cm]
    
    \includegraphics[width=0.25\linewidth]{uni_logo}\\[0.9cm]
    
    {\large\bfseries \curso}\\[0.6cm]
    
    {\LARGE\bfseries \tema}\\[0.35cm]
    {\normalsize\itshape \subtema}\\[1.4cm]
    
    \begin{minipage}{0.85\linewidth}
        \centering
        {\large\bfseries ALUMNO:}\\[0.2cm]
        {\normalsize \alumno \quad (\textbf{Código:} \codigo)}\\[0.7cm]
        
        {\large\bfseries DOCENTE:}\\[0.2cm]
        {\normalsize \docente}
    \end{minipage}
    
    \vfill
    {\normalsize \fecha}
    \vspace*{0.3cm}
\end{titlepage}

%% ------------------------------------------------------------
%%  DESARROLLO DEL TRABAJO (PÁGINA 1 DE CONTENIDO)
%% ------------------------------------------------------------
\newpage
\pagenumbering{arabic}
\setcounter{page}{1}

\section{Introducción: El estancamiento de la metafísica}

A finales del siglo XVIII, Immanuel Kant planteó una pregunta fundamental: \textit{¿Por qué disciplinas como la matemática y la física avanzan sólidamente construyendo acuerdos universales, mientras que la metafísica sigue estancada en las mismas disputas desde la antigüedad sin llegar a una sola verdad firme?}

En el Prólogo a la segunda edición (1787) de su obra cumbre, la \textit{Crítica de la razón pura}, Kant explica que una disciplina sólo es auténtica ciencia cuando entra en el «camino seguro de la ciencia». La matemática y la física lo lograron gracias a una revolución metodológica. En cambio, la metafísica tradicional quedó atrapada en un «mero andar a tientas» (\textit{Herumtappen}) y convertida en un campo de batalla (\textit{Kampfplatz}) de disputas interminables. Este informe analiza por qué la metafísica no puede llegar a ser una ciencia teórica demostrable.

\section{¿Cómo lograron la matemática y la física convertirse en ciencias?}

Para entender el estancamiento de la metafísica, Kant examina qué permitió el éxito de las disciplinas científicas:

\begin{itemize}
    \item \textbf{La matemática:} Se convirtió en ciencia cuando los geómetras griegos comprendieron que, para demostrar una propiedad geométrica (como en el triángulo isósceles), no debían medir figuras empíricas en la arena, sino \textbf{construir el concepto mediante reglas a priori} y deducir sus propiedades a partir de lo que la propia razón introdujo en él.
    
    \item \textbf{La física moderna:} Dio su salto definitivo cuando científicos como Galileo, Torricelli y Stahl abandonaron la observación pasiva. Galileo no esperó a ver cómo caían cuerpos al azar; construyó planos inclinados con pesos determinados para forzar a la naturaleza a responder ante condiciones controladas. Torricelli ideó un experimento en el que el aire debía sostener una columna de mercurio prevista por su hipótesis.
\end{itemize}

Kant sintetiza esto con una metáfora contundente: \textbf{la razón no debe acudir a la naturaleza como un alumno} pasivo que acepta todo lo que el maestro dice, sino como un \textbf{juez en ejercicio} que formula preguntas precisas y obliga a los testigos a responder de acuerdo con un plan previo. La ciencia progresa porque la mente propone principios teóricos y luego los somete a contraste empírico.

\section{La propuesta de Kant: El «Giro Copernicano»}

Al ver que la metafísica fracasaba al intentar conocer los objetos directamente, Kant propone un cambio de perspectiva análogo al de Nicolás Copérnico en astronomía: el célebre \textbf{giro copernicano}.

Antes de Copérnico, se intentaba explicar el cosmos asumiendo que la Tierra estaba fija en el centro y todos los astros giraban a su alrededor, lo que creaba contradicciones continuas. Copérnico ensayó la hipótesis inversa: supuso que el observador se mueve y los astros permanecen quietos, logrando un modelo astronómico ordenado y coherente.

Kant traslada esta lógica a la teoría del conocimiento:
\begin{itemize}
    \item \textbf{Enfoque tradicional (fallido):} Suponía que nuestra mente debía amoldarse pasivamente a los objetos para captarlos tal como son «en sí mismos». Bajo este esquema era imposible explicar cómo poseemos conocimientos universales y necesarios previos a la experiencia.
    \item \textbf{Revolución kantiana:} Debemos suponer que \textbf{los objetos de la experiencia se amoldan a nuestra forma humana de conocer}. No conocemos las cosas «en sí mismas» (\textit{noúmeno}), sino únicamente los fenómenos: es decir, tal como se nos presentan a través de nuestra sensibilidad (espacio y tiempo) y de las categorías lógicas del entendimiento (causa-efecto, sustancia, etc.).
\end{itemize}

\newpage

%% ------------------------------------------------------------
%%  DESARROLLO DEL TRABAJO (PÁGINA 2 DE CONTENIDO)
%% ------------------------------------------------------------

\section{Por qué la metafísica no puede llegar a ser una ciencia}

A partir del giro copernicano, la razón por la cual la metafísica fracasa como ciencia se vuelve evidente y definitiva:

\textbf{1. El requisito indispensable del conocimiento humano:} Para Kant, el conocimiento objetivo sólo surge cuando se unen de forma inseparable dos factores: las \textit{intuiciones sensibles} (datos concretos captados por los sentidos en el espacio y el tiempo) y los \textit{conceptos del entendimiento} (reglas mentales que ordenan y dan sentido a esos datos). Kant lo resume en su máxima clásica: \textit{«Los pensamientos sin contenido están vacíos; las intuiciones sin conceptos son ciegas»}. Un concepto abstracto sin un objeto real sensible que lo respalde es una cáscara vacía; e inversamente, una acumulación de percepciones sin conceptos no pasa de ser un caos sin sentido.

\textbf{2. El objeto propio de la metafísica (lo suprasensible):} La metafísica tradicional busca responder preguntas sobre tres grandes temas: el \textit{alma} (su naturaleza inmaterial e inmortal), el \textit{mundo como totalidad} (si tuvo inicio o es infinito, y si hay libertad o determinismo ciego) y \textit{Dios} (causa suprema de toda la realidad). Estos tres objetos son \textbf{suprasensibles}: están totalmente fuera del alcance de la experiencia sensible. Nadie puede ver el alma, medir el infinito cósmico con un instrumento ni observar empíricamente a Dios en un laboratorio.

\textbf{3. El fallo insalvable: conceptos vacíos y contradicciones insolubles:} Al carecer de intuición sensible, la metafísica comete el error de aplicar las categorías mentales fuera del ámbito de la experiencia posible. En consecuencia, opera con \textbf{conceptos totalmente vacíos de contenido empírico}. Como no existen observaciones posibles para contrastar sus hipótesis, no hay prueba objetiva ni experimento que permita verificar si una tesis metafísica es verdadera o falsa. Por ello, la razón cae en contradicciones lógicas insolubles (\textit{antinomias}): se puede demostrar con idéntica lógica que el universo tiene un comienzo en el tiempo y que es infinito, sin que la experiencia pueda desempatar jamás.

Por lo tanto, \textbf{la metafísica teórica nunca podrá ser una ciencia}, pues carece del requisito indispensable de todo saber científico: la posibilidad de contrastar sus afirmaciones con la experiencia sensible.

\section{Conclusión: El verdadero valor de la metafísica según Kant}

Que la metafísica no sea una ciencia no significa que deba desecharse. Kant sostiene que su valor auténtico es \textbf{crítico y ético}:
\begin{itemize}
    \item \textbf{Límite contra el dogmatismo:} Enseña a la ciencia a no extralimitarse pretendiendo saberlo todo de forma absoluta, frenando el dogmatismo y las supersticiones sin fundamento.
    \item \textbf{Fundamento de la moral:} Al probar que la razón teórica no puede demostrar a Dios ni al alma, Kant demuestra que \textbf{la ciencia materialista tampoco puede destruirlos ni refutarlos}. Así, la libertad humana queda preservada para la \textbf{razón práctica}: la ética, el deber moral y la responsabilidad del ser humano.
\end{itemize}

Como resumió el propio Kant: \textit{«Tuve que suprimir el saber [teórico] para dejar sitio a la fe [moral]»}. La metafísica no es una ciencia de hechos de la naturaleza, sino una brújula reguladora que delimita el conocimiento humano y da sentido a la vida moral.

\vspace{0.35cm}

\section*{Referencia bibliográfica}
\begin{itemize}
    \item Kant, I. (1787). \textit{Crítica de la razón pura} (Prólogo a la 2.ª ed., B VII--B XLIV). Trad. de Pedro Ribas (2006). Madrid: Taurus / Alfaguara.
\end{itemize}

\end{document}
